% For printing in a4
\documentclass[a4,10pt,twoside,openright,italian,english]{book}% twoside!

% For printing with the A5 format
%\documentclass[10pt,twoside,openright,english,italian]{book}% twoside!

% Set paper size
\usepackage[twoside=true]{geometry}

%For printing with the weird format
%\geometry{
%	paperwidth=17cm,
%	paperheight=24cm,
%	margin=2cm,
%	top=2.3cm,
%	bindingoffset=0.4cm
%}
% For printing in a4
\geometry{a4paper,
  margin=3cm,
  top=3.8cm,
  bindingoffset=0.4cm
}

%Uncomment this for final prints: this just enables printing on a4 paper
\usepackage[cam,center,a4,pdflatex,axes]{crop}

\usepackage{phdthesis}

\usepackage{fancyhdr}
\usepackage{color}
\usepackage{array}
\usepackage{mdwmath}
\usepackage{mdwtab}
\usepackage{amsmath,amssymb}
\usepackage{cite}
\usepackage{graphicx}
\usepackage{listings}
\usepackage{subfig}
\usepackage{booktabs}
\usepackage{latexsym}
\usepackage{color}
\usepackage{url}
\usepackage{bnf}
\usepackage{rotating}
\usepackage{multirow}
\usepackage{phdtitle}
\usepackage{paralist}
\usepackage{bibentry}
%\usepackage[algochapter]{algorithm2e}
\usepackage[bookmarks=true,
%pdftex=false,
bookmarksopen=true,hidelinks]{hyperref}
\usepackage[toc,acronym]{glossaries}
\usepackage{lscape}
\usepackage{algorithmic}
\usepackage{algorithm}
\usepackage{longtable}
\usepackage[T1]{fontenc} 
%\usepackage[latin1]{inputenc}
\usepackage[utf8]{inputenc}
%\usepackage{fontspec}
%\setmainfont{Calibri}
%% \hyphenation{} is used to force the 
\hyphenation{}

\newtheorem{Definition}{Definition}[section]

\lstset{tabsize=2,basicstyle=\footnotesize,breaklines=true}

\usepackage[english,italian]{babel}
\usepackage{ulem}
\normalem

\hypersetup{pdftitle={Title}, pdfauthor={Author}}

\nobibliography*

\department {Dottorato di ricerca in Ingegneria dell'Informazione}

% Please fulfil the followed fields with your data
\author{Fabio Carrara} 
\title{Title}
\tutor{Tutor Name}
\supervisor{Coordinator Name}
\titleimage{img/unipi.png} %please not change
\phdyear{Year}
\phdmonth{Month}
\phdcycle{Cycle XXXII ?}

%%%%%%%%%%%%%%%%%%%%%%%%%%%%%%%%%%%%%%
% Let's Start The Real Document
%%%%%%%%%%%%%%%%%%%%%%%%%%%%%%%%%%%%%%
\newglossaryentry{matrix_channel}
{       name={$H^*$},
        description={Conjugate operation}
}
\newglossaryentry{trasp_x}
{       name={$[ x ]^{\rm T}$},
        description={transpose operator}
}
\newglossaryentry{vec_x}
{       name={\textbf{x}},
        description={vectors are in bold}
}
\newglossaryentry{floor_funct}
{       name={$\left\lfloor x \right\rfloor$},
        description={round to the lower integer of $x$}
}

% Deep Learning
\newacronym{ml}{ML}{Machine Learning}
\newacronym{dl}{DL}{Deep Learning}
\newacronym{dnn}{DNN}{Deep Neural Network}
\newacronym{ffnn}{FFNN}{Feed-Forward Neural Network}
\newacronym{rnn}{RNN}{Recurrent Neural Network}
\newacronym{mlp}{MLP}{Multilayer Perceptron}
\newacronym{cnn}{CNN}{Convolutional Neural Network}
\newacronym{dcnn}{DCNN}{Deep Convolutional Neural Network}
\newacronym{lstm}{LSTM}{Long Short-term Memory}
\newacronym{sgd}{SGD}{Stochastic Gradient Descent}
\newacronym{adam}{Adam}{Adaptive Moment Estimation}
\newacronym{relu}{ReLU}{Rectified Linear Unit}
\newacronym{resnet}{ResNet}{Residual Network}
\newacronym{senet}{SENet}{Squeeze-and-Excitation Network}
\newacronym{bn}{BN}{Batch Normalization}
\newacronym{rmac}{R-MAC}{Regional Maximum Activation of Convolutions}

% Metrics
\newacronym{tpr}{TPR}{True Positive Rate}
\newacronym{tnr}{TNR}{True Negative Rate}
\newacronym{fpr}{FPR}{False Positive Rate}
\newacronym{fnr}{FNR}{False Negative Rate}
\newacronym{auc}{AUC}{Area Under the Curve}
\newacronym{roc}{ROC}{Receiver Operation Characteristics}
\newacronym{eer}{EER}{Equal Error Rate}
\newacronym{ap}{AP}{Average Precision}
\newacronym{map}{mAP}{mean Average Precision}
\newacronym{dcg}{DCG}{Discounted Cumulative Gain}
\newacronym{ndcg}{NDCG}{Normalized Discounted Cumulative Gain}
\newacronym{r@k}{R@K}{Recall at K}
\newacronym{medR}{medR}{Median Rank}
\newacronym{mrr}{MRR}{Mean Reciprocal Rank}

% CBIR
\newacronym{cbir}{CBIR}{Content-based Image Retrieval}
\newacronym{pca}{PCA}{Principal Component Analysis}

% Datasets
\newacronym{ilsvrc}{ILSVRC}{ImageNet Large Scale Visual Recognition Challenge}
\newacronym{coco}{COCO}{Common Objects in Context} % TODO check COCO acronym
\newacronym{t4sa}{T4SA}{Twitter for Sentiment Analysis}

% Others
\newacronym{gpu}{GPU}{Graphical Processing Unit}
\newacronym{lbp}{LBP}{Local Binary Patterns}
\newacronym{lpq}{LPQ}{Local Phase Quantization} % CHECK
\newacronym{str}{STR}{Surrogate Text Representation}
\newacronym{svm}{SVM}{Support Vector Machine}

\newacronym{vgg}{VGG}{Visual Geometry Group}



\makeglossaries

\begin{document}
\selectlanguage{english}

\maketitle

\pagestyle{empty}

\cleardoublepage
\newpage

%%%%%%%%%%%%%%%%%%%%%%%%%%%%%%%%%%%%%%
% Dedication - For removing dedication, please comment or delete the code inside \thispagestyle{empty}
%%%%%%%%%%%%%%%%%%%%%%%%%%%%%%%%%%%%%%
\thispagestyle{empty}
    \null\vspace{\stretch {1}}
        \begin{flushright}
                This thesis is dedicated to....
        \end{flushright}
\vspace{\stretch{2}}\null

\cleardoublepage
\newpage

\pagestyle{empty}
%% change numbering into Roman numbers for the introductory part
%\setcounter{page}{1}
%\pagenumbering{Roman}

%%%%%%%%%%%%%%%%%%%%%%%%%%%%%%%%%%%%%%
% Quotes - For removing quotes, please comment or delete the code inside \thispagestyle{empty}
%%%%%%%%%%%%%%%%%%%%%%%%%%%%%%%%%%%%%%
\thispagestyle{empty}
    \null\vspace{\stretch {1}}
        \begin{flushright}
                "A few first rate research papers are preferable to a large number \\
                that are poorly conceived or half-finished.\\
                The latter are no credit to their writers and \\
                a waste of time to their readers"\\
                Claude Shannon
                %IRE Transactions on Information Theory (1956), volume 2, issue 1, page 3. Shannon, Claude E. (March 1956), The Bandwagon, 2, doi:10.1109/TIT.1956.1056774.
        \end{flushright}
\vspace{\stretch{2}}\null

\cleardoublepage
\newpage

\pagestyle{empty}
%% change numbering into Roman numbers for the introductory part
\setcounter{page}{1}
\pagenumbering{Roman}

%%%%%%%%%%%%%%%%%%%%%%%%%%%%%%%%%%%%%%
% Acknowledgement
%%%%%%%%%%%%%%%%%%%%%%%%%%%%%%%%%%%%%%
\chapter*{Acknowledgements}
\lettrine{A}{cknowledgements} goes here.

\selectlanguage{italian}
%\chapter*{Ringraziamenti}
%\lettrine{R}{ingraziamenti}
% English-only Acknowledgements
\selectlanguage{english}

\cleardoublepage
\newpage

\pagestyle{fancy}
% change numbering into Roman numbers for the introductory part

%%%%%%%%%%%%%%%%%%%%%%%%%%%%%%%%%%%%%%
% Summary
%%%%%%%%%%%%%%%%%%%%%%%%%%%%%%%%%%%%%%
\selectlanguage{english}
\chapter*{Summary}
\lettrine{T}{he} large diffusion of cheap cameras and smartphones led to an exponential daily production of digital visual data, such as images and videos.
In this context, most of the produced data lack manually assigned metadata needed for their manageability in large-scale scenarios, thus shifting the attention to the automatic understanding of the visual content.
Recent developments in Computer Vision and Artificial Intelligence empowered machines with high-level vision perception enabling the automatic extraction of high-quality information from raw visual data.
Specifically, \acrfullpl{cnn} provided a way to automatically learn effective representations of images and other visual data showing impressive results in vision-based tasks, such as image recognition and retrieval.

In this thesis, we investigated and enhanced the usability of \acrshortpl{cnn} for visual data management.
First, we identify three main limitations encountered in the adoption of \acrshortpl{cnn} and propose general solutions that we experimentally evaluated in the context of image classification.
We proposed miniaturized architectures to decrease the usually high computational cost of \acrshortpl{cnn} and enable edge inference in low-powered embedded devices.
%We evaluated our proposal in a practical distributed application, i.e., visual parking lot occupancy detection, extensively comparing the reduced models to state-of-the-art methods and architectures.
We tackled the problem of manually building huge training sets for models by proposing an automatic pipeline for training classifiers based on cross-media learning and Web-scraped weakly-labeled data.
We analyzed the robustness of \acrshortpl{cnn} representations to out-of-distribution data, specifically the vulnerability to adversarial examples, and proposed a detection method to discard spurious classifications provided by the model.
%
Secondly, we focused on the integration of \acrshort{cnn}-based \acrfull{cbir} in the most commonly adopted search paradigm, that is, textual search.
We investigated solutions to bridge the gap between image search and highly-developed textual search technologies by reusing both the front-end (text-based queries) and the back-end (distributed and scalable inverted indexes). % in \acrshort{cbir} scenarios.
We proposed a cross-modal image retrieval approach which enables textual-based image search on unlabeled collections by learning a mapping from textual to high-level visual representations.
Finally, we formalized, improved, and proposed novel surrogate text representations, i.e., text transcriptions of visual representations that can be indexed and retrieved by available textual search engines enabling \acrshort{cbir} without specialized indexes.

\selectlanguage{english}

%%%%%%%%%%%%%%%%%%%%%%%%%%%%%%%%%%%%%%
% Italian Summary
%%%%%%%%%%%%%%%%%%%%%%%%%%%%%%%%%%%%%%

%There must be an Italian version of the summary
\selectlanguage{italian}
\chapter*{Sommario}
\lettrine{L'}{} enorme diffusione di fotocamere e smartphone a prezzi economici ha portato a una produzione esponenziale giornaliera di dati visivi digitali, come immagini e video.
La maggior parte dei dati prodotti non è accompagnata dai metadati, come descrizioni, tags, o altri dati manualmente assegnati, che sono necessari per la loro gestione automatica su larga scala.
L'attenzione della ricerca si è quindi spostata sulla comprensione automatica del contenuto visivo di tali dati.
I recenti sviluppi dell'Intelligenza Artificiale applicata alla Computer Vision hanno reso possibile l'estrazione automatica di informazioni di alta qualità direttamente da dati visivi grezzi (pixel).
In particolare, modelli neurali come le reti neurali convoluzionali hanno fornito un modo per apprendere automaticamente delle rappresentazioni numeriche efficaci per immagini e altri dati visivi che hanno ottenuto risultati impressionanti in task visivi come il riconoscimento di immagini.

In questa tesi, è stata studiata e migliorata l'usabilità delle reti neurali convoluzionali per la gestione dei dati visivi su larga scala.
Nella prima parte, sono state identificate tre limitazioni principali solitamente incontrate nell'utilizzo delle reti convoluzionali e sono state proposte delle soluzioni generali che abbiamo valutato sperimentalmente nel contesto della classificazione di immagini.
Sono state proposte architetture miniaturizzate per ridurre il costo computazionale solitamente elevato di questo tipo di reti e consentire quindi il loro utilizzo anche a bordo di dispositivi embedded a bassa potenza.
% Abbiamo valutato la nostra proposta in un'applicazione pratica e distribuita, ovvero la rilevazione di occupazione del parcheggio visivo, confrontando ampiamente i modelli ridotti con i metodi e le architetture più avanzati.
È stato affrontato il problema della creazione di training set per i modelli, che richiederebbero un notevole sforzo manuale, proponendo una pipeline automatica per allenare reti basata sul cross-media learning e su dati imprecisi provenienti dal Web.
È stata analizzata la robustezza delle rappresentazioni estratte dalle reti convoluzionali per dati fuori dalla distribuzione di train, con enfasi particolare sulla vulnerabilità delle reti agli attacchi avversari (adversarial examples), proponendo un metodo di rilevamento per scartare le classificazioni spurie fornite dal modello.
%
In secondo luogo, ci siamo concentrati sull'integrazione della ricerca per immagini, basata su rappresentazioni estratte da reti convoluzionali, col paradigma di ricerca più comunemente adottato, cioè la ricerca testuale.
In questo contesto, abbiamo studiato delle soluzioni per colmare il divario tra l'attuale stato dell'arte nella ricerca di immagini e le più mature tecnologie di ricerca testuale.
In particolare, sono state integrate soluzioni per la ricerca di immagini basata sul contenuto sia con il front-end (query testuali) che con il back-end (indici invertiti distribuiti e scalabili per documenti testuali).
Nel primo caso, è stato proposto un approccio di recupero di immagini cross-modale che consente la ricerca tramite descrizione testuale di immagini in collezioni non etichettate tramite l'apprendimento di una funzione di mapping delle rappresentazioni testuali in quelle visive.
Nel secondo caso, sono state formalizzate, migliorate e proposte nuove rappresentazioni testuali surrogate per immagini, che consistono in una trasformazione delle rappresentazioni visive in testo surrogato che può essere indicizzato e recuperato dai motori di ricerca testuali attualmente disponibili, abilitando applicazioni di recupero di immagini senza il bisogno di indici specializzati.

\selectlanguage{english}

%%%%%%%%%%%%%%%%%%%%%%%%%%%%%%%%%%%%%%
% Publications
%%%%%%%%%%%%%%%%%%%%%%%%%%%%%%%%%%%%%%

%List of publications of the PhD candidate
\selectlanguage{english}
\chapter*{List of publications}

% for citing your paper, take directly APA format from Google scholar
% \item Surname, N., Surname, N. and Surname, N. (Year,Month). Title of the paper or journal. \emph{Place of publication}. (Vol. 123, pp. 456). EditorName.

\section*{International Journals}
\begin{enumerate}
    \item Amato, G., Carrara, F., Falchi, F., Gennaro, C., Meghini, C., and Vairo, C. (2017, April). Deep learning for decentralized parking lot occupancy detection. \emph{Expert Systems with Applications}. (Vol. 72, pp. 327-334). Pergamon.
    \item Carrara, F., Esuli, A., Fagni, T., Falchi, F., and Fernández, A. M. (2018, June). Picture it in your mind: Generating high level visual representations from textual descriptions. \emph{Information Retrieval Journal}. (Vol. 21, pp. 208-229). Springer.
    % manca mese, volume e pagine
    \item Carrara, F., Falchi, F., Caldelli, R., Amato, G., and Becarelli, R. (2018). Adversarial image detection in deep neural networks. \emph{Multimedia Tools and Applications}. (pp. 1-21). Springer.

\end{enumerate}

\section*{International Conferences/Workshops with Peer Review}
\begin{enumerate}
    \item Amato, G., Carrara, F., Falchi, F., Gennaro, C., and Vairo, C. (2016, June). Car parking occupancy detection using smart camera networks and deep learning. In \emph{2016 IEEE Symposium on Computers and Communication (ISCC)}. (pp. 1212-1217). IEEE.
    \item Carrara, F., Falchi, F., Caldelli, R., Amato, G., Fumarola, R., and Becarelli, R. (2017, June). Detecting adversarial example attacks to deep neural networks. In \emph{Proceedings of the 15th International Workshop on Content-Based Multimedia Indexing (CBMI)}. (p. 38). ACM.
    \item Amato, G., Carrara, F., Falchi, F., and Gennaro, C. (2017, June). Efficient Indexing of Regional Maximum Activations of Convolutions using Full-Text Search Engines. In \emph{Proceedings of the 2017 ACM on International Conference on Multimedia Retrieval (ICMR)}. (pp. 420-423). ACM.
    \item Vadicamo, L., Carrara, F., Cimino, A., Cresci, S., Dell’Orletta, F., Falchi, F., and Tesconi, M. (2017, October). Cross-Media Learning for Image Sentiment Analysis in the Wild. In \emph{2017 IEEE International Conference on Computer Vision (ICCV) Workshops}. (pp. 308-317).
    \item Amato, G., Bolettieri, P., Carrara, F., Falchi, F., Gennaro, C. (2018, June). Large-Scale Image Retrieval with Elasticsearch. In \emph{Proceedings of the 41st International ACM Conference on Research and Development in Information Retrieval (SIGIR)}. (pp.  925-928). ACM.
\end{enumerate}

%\section*{Others}
%\begin{enumerate}
%    % for citing your paper, take directly APA format from Google scholar
%    \item Surname, N., Surname, N. and Surname, N. (Year,Month). Title of the paper or journal. \emph{Place of publication}. (Vol. 123, pp. 456). EditorName.
%\end{enumerate}

\selectlanguage{english}

%%%%%%%%%%%%%%%%%%%%%%%%%%%%%%%%%%%%%%
% List of Abbreviation and Symbols
%%%%%%%%%%%%%%%%%%%%%%%%%%%%%%%%%%%%%%
%Have a look at \gls command... you can refer a term in multiple ways
\printglossary[type=\acronymtype,title=List of Abbreviations]
\let\cleardoublepage\clearpage
\printglossary[title=Notation]
\cleardoublepage
\newpage
%%%%%%%%%%%%%%%%%%%%%%%%%%%%%%%%%%%%%%
% TOC
%%%%%%%%%%%%%%%%%%%%%%%%%%%%%%%%%%%%%%
\tableofcontents
\cleardoublepage
\newpage

% Now lets go back to normal numbering
\setcounter{page}{1}
\pagenumbering{arabic}

\cleardoublepage
\chapter{Title chapter 1}

%============================= INTRODUCTION =================================
\section{Introduction}\label{sec:i}
Remember that wherever you want you can cite with \verb.\cite{a1}. \cite{a1}, \cite{b1} and/or \cite{m1} someone.
Or use \verb.\footnote{That's a footnote}. like this\footnote{That's a footnote} 
% this is used to include other .tex files in order to split a long document 

%============================= SECTIONS =================================

\section{New section}
You can insert a definition
	\begin{Definition}
		Unipi: Universit\'a di Pisa
	\end{Definition}
	
\section{Another section}
	\subsection{And a subsection}
		In this sub section I will include many images just to make the list of figures meaningful.

			\begin{figure}[h!tb]
				\centerline {\includegraphics[scale=0.6]{img/unipi.png}}
				\caption{Caption of this Unipi image.}
				\label{fig:leet}
			\end{figure}

\section{New section}
	\subsection{And also many sub}
		\subsubsection{Or again more nested}
			\begin{itemize}
				\item item 1
				\item item 2
				\item item 3
			\end{itemize}			
			
\section{New section}
This is an example of table. You can see it in Table~\ref{tab:TableName}. 

 \begin{table}[h!tb]
   \centering \caption{Table caption}
   \label{tab:TableName}
   \vskip 0.2cm
   %%
   \scalebox{0.90}{
	    %% The {|c|c|c|c|c|} define the number of columns.
	    %% c means centered
	    %% | defines a vertical line between two columns 
	    \begin{tabular}{|c|c|c|}
	      \hline
	      Col 1 & Col 2 & Col2  \\
	      %% \\ force a newline without creating an horizontal line
	      Dim C.1 & Dim C.3 & Dim C.3 \\
	      %% \hline create an horizontal line between two rows
	      \hline
	      Data 1.1 & Data 1.2 & Data 1.3 \\
	      \hline
	      Data 2.1 & Data 2.2 & Data 2.3 \\
	      \hline
	      Data 3.1 & Data 3.2 & Data 3.3 \\  
	      \hline
	    \end{tabular}
	 }
 \end{table}

\section{You can do as many section as you want}
The following equation has been created using \verb.$$.: $ y_1 = a*x + b$.

%% \smallskip introduces a small vertical skip
\smallskip

%% \noindent is used to eliminate the blankline and the paragraph settings 
\noindent
This is a more complex equation: \[ y_2 = a*x + b \], created using \verb.\[ \].

%% \bigskip introduces a big vertical skip
\bigskip

\noindent
This is, (\ref{eq:eeq1}), an enumerated equation: 

	\begin{equation}\label{eq:eeq1}
		y_3 = a*x + b
	\end{equation}

, created using:

	\begin{verbatim}
		\begin{equation}\label{eq:eeq1}
		    y_3 = a*x + b
		\end{equation}
	\end{verbatim}

It is also possible to insert acronyms, after definition in \textit{Acronyms} file, like \acrshort{fpsLabel} which stands for \acrlong{fpsLabel} or \acrshort{SNRLabel} which stands for \acrlong{SNRLabel}.
Please define the notation at the beginning, defining operators and symbols in \textit{Notation} file, such as \gls{matrix_channel}, \gls{trasp_x}, \gls{vec_x}, \gls{floor_funct}. % uncomment if you want part I
% Insert further chapters using the following syntax
% \input{Source_Folder_Name/FileName}


\cleardoublepage
\phantomsection
\addcontentsline{toc}{chapter}{\bibname}
\small
\bibliographystyle{plain}
\bibliography{thesis}

\end{document}
